\chapter{EEG Signal}


\section{Overview}

Electroencephalography (EEG) is a non-invasive technique used to record electrical activity of the brain\cite{Mushtaq2024}.
It involves placing electrodes on the scalp to record electrical activity of the brain --- the electroencephalogram (also EEG).
Each electrode records the voltage fluctuations resulting from the synchronized discharges of the working neurons in the brain.
By subtracting the reference electrode reading from a reading of an electrode of interest we cancel out the common environmental noise and get a signal that constitutes a single channel of an EEG recording.
The amplitude of the resulting signal is in the order of microvolts and very sensitive to any head movements, which need to be minimized during a recording session.
Modern equipment captures these fluctuations with upwards of 1000 Hz temporal resolution, allowing for precise linking between the brain's recorded activity and external stimuli.


\section{History}

A British physician, Richard Caton, was the first to discover in 1875 that the brain --- specifically of rabbits and monkeys --- elicits electrical activity and, furthermore, that this activity seems related to its function\cite{caton}. For his discovery, he used a galvanometer, invented earlier in 1820 and used in laboratories and telecommunication.
At the time of his discovery, it was already known that nerves conduct electrical signals and that these signals are used to trigger muscle tissue to contract\cite{mattucci1}\cite{matteucci2}.
In 1901, a Dutch medical doctor, Willem Einthoven, invented a more precise string galvanometer and used it to record an electrocardiogram (ECG)\cite{about_einthoven}. By placing electrodes on the left and right arms and a leg, he measured the electrical signal produced by a beating heart. The device functioned by placing a fine conducting string within a magnetic field created by two powerful electromagnets; the electrical signal caused this string to deflect. He was later awarded the Nobel Prize in Physiology or Medicine for his discovery.
Later, in 1912, a Ukrainian physiologist, Vladimir Pravdich-Neminsky, used this invention and published the first EEG recording recorded from a dog\cite{about_vladimir}.

In 1924 --- which coincidently was the year when Einthoven was awarded his Nobel Prize --- a German psychiatrist, Hans Berger, while looking for the physiological basis for telepathy, recorded the first human EEG and discovered the ``alpha wave''—the frequency present in the signal when a person is awake\cite{about_berger}.
His discovery was initially met with skepticism\cite{Mushtaq2024}, which changed a decade later when Frederick Gibbs and others used EEG to describe specific, repeatable patterns in electrical activity during epilepsy\cite{about_gibson}.
Later, in the 50s, the EEG signal was used to research sleep.
In 1953, researchers discovered the REM (Rapid Eye Movement) sleep phase, characterized by rapid eye movements and high-frequency, low-voltage EEG activity, akin to alertness, which is often the phase when dreaming occurs\cite{rem_sleep}.
Later, a standardized way to divide sleep into sleep stages based on EEG was developed\cite{sleep_stage}.

Already in 1973, the first demonstration of a brain-computer interface (BCI) was seen, where the EEG signal was used to control a cursor on a screen\cite{vidal}.
This decade also saw the first use of computers to analyze the EEG signal, and later the analysis and collection of EEG data became purely digital.
In the 1980s, laboratories started syncing EEG recordings with video recordings of the subject, allowing for more detailed studies of the relationship between brain activity and behavior\cite{gloor}.

Modern-day 128- or 256-electrode systems allow for spatially detailed recordings of brain activity, including source localization.
EEG is now widely used in clinical settings for the diagnosis of sleep disorders\cite{sleep_medicine} and the localization and classification of seizures\cite{fisher}.
We also see a rise in the availability of consumer-grade portable EEG devices, which are used for various applications, including neurofeedback, meditation, and sleep improvements.

\section{Characteristics}

The EEG signal has a few characteristics that make it extra challenging to process and analyze.

\subsection{Non-stationarity}

A non-stationary process is one whose statistical properties --- like mean, variance, or autocorrelation --- change over time. 
The EEG signal is non-stationary\cite{nonstationary} because the brain's activity is dynamic and influenced by various internal and external factors, such as cognitive processes, sensory inputs, emotional states, and physiological changes like variations in skin impedance.
Without any adaptations, this renders many classical statistical methods of data analysis ineffective.
This suggests that the normalization of signal data should be done on a per-window basis.

\subsection{Artifacts}

Because the amplitude of neural oscillations as seen in the EEG signal is so low, requiring precise measurement, even small movements and external factors can introduce significant noise and artifacts into the recordings.
Specifically, any head movements or eye blinks cause large artifacts in the signal. It can also happen that an electrode loses contact with the skin or that heartbeats appear from an electrode placed near a large blood vessel.
Given the small sample size, this often necessitates marking and removing the artifacts from the data.

\subsection{Expensive and time-consuming data collection}

Collecting EEG data requires specialized, expensive equipment and qualified personnel.
It also requires finding subjects willing to participate in the study and fitting the study requirements.
For these reasons, EEG datasets are often small, counting hundreds to thousands of samples collected from up to a few dozen subjects.

\subsection{Noise}

Given the low amplitude of the EEG signal, it is very susceptible to any external noise.
Commonly, EEG data preprocessing includes filtering the 50 or 60 Hz power line noise, which appears as thick fuzz on the signal.
In studies that combine music listening with EEG collection, the music itself would be a source of noise.
To mitigate this, studies can use electromagnetically shielded insert earphones\cite{nozaradan}, which pump the sound directly into the ear via hollow plastic tubes, keeping the transducer (the part that creates sound) further away from the EEG apparatus.

\subsection{Signal-to-noise ratio}

Given that EEG measures a combined electrical activity of millions of neurons from a brain that reacts to and processes a multitude of internal and external stimuli, 
it is a given that the observed reaction to the experimental stimulus as witnessed from the data is going to be mixed with a lot of other brain activity, which is not related to the stimulus.
Furthermore, the observed reaction is going to be different across subjects and even across different recording sessions of the same subject.
For this reason, the signal-to-noise ratio (SNR) of the EEG data is often very low, making it difficult to extract meaningful information from the data and to link it to the experimental stimulus.
A common way of analyzing EEG data involves averaging the EEG response at fixed intervals after presenting the stimulus.

\section{Preprocessing}

To mitigate the above challenges, various preprocessing techniques are used to clean the data and extract relevant features.
% Some of these techniques became standardized in a form of python or matlab libraries, which are widely used in the field.
We generally consider the EEG signal to be a two-dimensional array of shape (n\_channels, n\_timepoints), where n\_channels is the number of electrodes used in the recording and n\_timepoints is the number of time points recorded for each channel.

\subsection{Filtering}

Filtering is a crucial preprocessing step in EEG signal processing. It involves removing unwanted frequency components from the signal, which can be caused by various sources of noise and artifacts.
Commonly, the signal is processed through a high-pass filter (e.g., 1Hz) to remove slow drifts caused by factors like sweating and changes in skin impedance.
It can also be processed through a low-pass filter (e.g., 30Hz-100Hz) to remove high-frequency noise and focus on much better-understood lower frequencies in the EEG data.
The power noise is removed with a notch filter at 50 or 60 Hz, depending on the local power line frequency of the recording environment.

\subsubsection{Band power processing}

A common signal processing technique --- also applied to EEG data --- is to filter the signal into chosen frequency bands and then calculate the power of the signal in each band, which is a measure of the strength of the oscillations in that band.
Over the years, an array of research has been done linking specific frequency bands in the EEG signal to specific cognitive processes, with different roles linked to different frequency bands\cite{basar}.
This guides the choice of frequency bands to be used for band power processing.

\begin{table}[h]
\centering
\begin{tabular}{lll}
\hline
\textbf{Band} & \textbf{Frequency range} & \textbf{Associated states / processes} \\
\hline
Delta   & 0.5--4 Hz  & Deep sleep, unconscious processes\cite{harmony} \\
Theta   & 4--8 Hz    & Drowsiness\cite{lal}, meditation\cite{lagapulos}, memory encoding\cite{klimesch} \\
Alpha   & 8--13 Hz   & Relaxed wakefulness, eyes closed, idle state\cite{alastair} \\
Beta    & 13--30 Hz  & Active thinking, concentration, alertness\cite{engel}\cite{ray} \\
Gamma   & 30--100 Hz & Higher cognitive functions\cite{bertrand}\cite{fries} \\
\hline
\end{tabular}
\caption{EEG frequency bands and their associated cognitive states.}
\label{tab:eeg_bands}
\end{table}

\subsection{Independent Component Analysis}

Independent Component Analysis (ICA) is a method of separating recorded EEG channels into independent components.
Concretely, naming the EEG signal \( \mathbf{x} \) and the result of applying ICA \( \mathbf{s} \), ICA finds an inverse linear transformation to a matrix \( \mathbf{A} \) and some \( \mathbf{s} \):

\[
\mathbf{x} = \mathbf{A}\mathbf{s}
\]

such that the channels of \( \mathbf{s} \) are:
\begin{itemize}
  \item{statistically independent}
  \item{non-Gaussian}
  \item{linearly mixed in the observed signal \( \mathbf{x} \)}
\end{itemize}
Here \( \mathbf{x} \) is of dimension (n\_channels, n\_timepoints) 
and \( \mathbf{s} \) is of dimension (n\_components, n\_timepoints), 
where n\_components is chosen and \( \mathbf{A} \) is of dimension (n\_channels, n\_components).

In this formulation, the solution \( \mathbf{s} \) is not unique, and specifically, there is no inherent ordering to the independent components.
Also, the transform is applied independently to each recorded EEG snippet, making the subsequent analysis a harder problem.
To mitigate this, an ordering can be fixed by sorting the independent components by how much variance of the original signal they explain (Variance Accounted For (VAF)).

\subsection{Artifact classification and removal}

Once a signal is decomposed into independent components, we gain a potent tool for artifact removal.
By classifying a component into what type of activity it represents --- e.g., eye blinks, heartbeats, muscle activity, brain activity --- we can choose to remove the components that are classified as artifacts and reconstruct the signal without them.
The classification can be done with a neural network trained on expert-labeled data\cite{iclabel}. Modern-day trained classifiers like this are a common inclusion in EEG processing libraries.

\subsection{Normalization}

There are a few reasonable ways to normalize the signal, such as z-score normalization, min-max normalization, or robust normalization.
Because of the non-stationarity of the signal, it is common to apply normalization on a per-window basis, where a window is a short segment of the full recording.
What remains to be chosen is whether the normalization should apply per channel independently or jointly across all channels.
Simply removing the mean across all channels is called rereferencing.


% historia aparatury:
%  - pierwsze aparaty
%  - aparaty współczesne
%  - ! system 10-20:
%   * bcmi 10-20
%   * musing: GSN-HydroCel-129 https://gemini.google.com/app/7785213939c45409

% charakterystyki:
% - niestacjonarne
% - szum
% - artefakty

% preprocessing:
%  - filtracja
%  - ICA
%  - artifact classification and removal
%  - band power
%  - averaging
